% ============================================================================
\chapter{PHP : les bases, le langage serveur le plus direct qui soit}
\label{ch:php-bases}
% ============================================================================

\begin{objectif}
\begin{itemize}
  \item Comprendre le cycle d'une requ\^ete PHP et l'installer en local en trois commandes.
  \item \'Ecrire du PHP moderne~: typage strict, tableaux, fonctions, fermetures.
  \item Ma\^itriser le type \og ~array~ \fg{}, qui est en r\'ealit\'e une table ordonn\'ee, et ses fonctions.
  \item Structurer un projet PHP sans framework~: un front controller, des fonctions, de la clart\'e.
  \item Conna\^itre les erreurs de d\'ebutant (headers, encodage, \verb|==|, includes) et les \'eviter.
\end{itemize}
\end{objectif}

\section{PHP, c'est quoi exactement}

PHP est un moteur qui ex\'ecute un script \emph{\`a chaque requ\^ete}, puis meurt. Il n'y a pas d'\'etat
persistant entre deux requ\^etes~: tout ce que tu ne stockes pas (base de donn\'ees, fichier, session,
cache m\'emoire) est perdu quand le script se termine. Cette contrainte explique \`a la fois la
simplicit\'e de PHP (pas de serveur \`a garder vivant, pas d'\'ev\'enements, pas d'async) et ses pi\`eges.

\begin{enumerate}
  \item Une requ\^ete HTTP arrive (Apache avec \verb|mod\_php|, ou nginx + \verb|php-fpm|).
  \item le moteur \textbf{compile} le script en opcode, le met en cache (\textbf{OPcache})~;
  \item le script s'ex\'ecute de la premi\`ere \`a la derni\`ere ligne~; il produit du texte (souvent du HTML)~;
  \item la r\'eponse est envoy\'ee, la m\'emoire est lib\'er\'ee. Point final.
\end{enumerate}

Version de r\'ef\'erence pour ce livre~: \textbf{PHP 8.3 / 8.4}. Tout ce qui est marqu\'e \og ~PHP~8+~ \fg{}
n'existe pas en 7.4 (finie de support, encore trop fr\'equente sur les vieilles mutualisations).

\begin{terminal}
# Verifier / installer
php -v
# Debian/Ubuntu : sudo apt install php-cli php-fpm php-mysql php-curl php-mbstring php-xml php-zip
# macOS          : brew install php
# Windows        : https://windows.php.net/download/ (voir chap. 1)
# Alternative : conteneur
docker run --rm -it -v $PWD:/app -w /app php:8.3-cli php script.php

# Le serveur de developpement integre : remplace apache pour apprendre
php -S localhost:8000 -t public
php -S 127.0.0.1:8000 public/router.php       # avec un routeur pour simuler les rewrites
\end{terminal}

\begin{note}
Le serveur int\'egr\'e (\verb|php -S|) est mono-thread~: une seule requ\^ete \`a la fois, parfait pour
apprendre, jamais pour un test de charge. Il refuse aussi les connexions entrantes autres que
localhost par d\'efaut~: utilise \verb|0.0.0.0:8000| pour tester depuis un t\'el\'ephone.
\end{note}

\section{Premier script, et les r\`egles de syntaxe}

\begin{codephp}
<?php
declare(strict_types=1);                    // PRECIEUX : interdit la conversion implicite des types
header('Content-Type: text/html; charset=utf-8');

$prenom = 'Nour';
$age = 28;
printf('%s a %d ans.%s', $prenom, $age, PHP_EOL);

// Une balise ouvrante, une seule par fichier ; JAMAIS de balise fermante dans un fichier
// qui ne contient que du PHP (sinon tu ajoutes un saut de ligne a la reponse, et les
// en-tetes "headers already sent" deviennent un cauchemar).
\end{codephp}

\begin{tabz}{lX}
  \headrow \thd{R\`egle} \& \thd{D\'etail qui compte} \\
  \midrule
  \verb|;| obligatoire \& chaque instruction se termine par un point-virgule. Le \verb|?>| final en implique un. \\
  \verb|\$| devant toute variable \& \verb|\$x|~; \verb|x| est une constante ou un nom de fonction. \\
  Insensible \`a la casse pour les mots-cl\'es/fonctions \& \verb|ECHO|, \verb|strlen| et \verb|StrLen| sont \'egaux. \textbf{Les noms de classes le sont aussi}~: d'o\`u les conventions PSR. \\
  Sensible pour les variables \& \verb|\$nom| et \verb|\$Nom| sont deux variables diff\'erentes. \\
  Cha\^ines \& \verb|'simple'| = litt\'erale~; \verb|"double"| = interpol\'ee (\verb|"\$a[0]|, \verb|"{\$obj->p}|). \\
  Commentaires \& \verb|//|, \verb|\#|, \verb|/* */|, et \verb|/** DocBlock */| pour la documentation lue par les IDE. \\
  \verb|echo| / \verb|print| / \verb|printf| \& \verb|echo| plusieurs valeurs~; \verb|print| renvoie 1~; \verb|printf| formate. \\
\end{tabz}

\section{Types, variables, et le typage moderne}

\begin{codephp}
<?php
declare(strict_types=1);

// Types scalaires
$entier = 42;            // int (64 bits, sauf sur 32 bits : 2^31-1 -- un id de tweet qui casse, ca existe)
$flottant = 3.14;        // float ; 1e3, 0x1A, 0b1010, 1_000_000 sont valides
$chaine = "cinq";        // string : une suite d'octets, pas de notion de caractere unicode
$booleen = true;         // bool
$rien = null;            // null : "absence de valeur"
$tableau = [1, 2, 3];    // array
$objet = new stdClass(); // object
$ressource = fopen('php://memory', 'r');   // ressource (fermeture explicite recommandee)

// Declarations typ\'ees : la norme en PHP moderne
function compter(array $articles, ?string $filtre = null, int $limite = 20): int
{
    $n = 0;
    foreach ($articles as $a) {
        if ($filtre === null || str_contains((string) $a['titre'], $filtre)) {
            $n++;
            if ($n >= $limite) { break; }
        }
    }
    return $n;
}

// Types d'union et d'intersection, nommes, readonly (8.1+), jamais (8.1)
function id(): int|string { return random_int(1, 1000); }
function traiter(Countable&ArrayAccess $c): never { throw new RuntimeException('non'); }

// Propriete typ\'ee : pas de valeur par defaut implicite -> une erreur nette si non initialisee
final class Point { public function __construct(public float $x = 0.0, public readonly float $y = 0.0) {} }

// gettype() est un outil de debug ; get_debug_type() (8.0) est plus lisible
// var_dump() : le reflexe. settype(), intval() : a eviter, prefere les casts (int), (float), (string), (bool), (array)

// Comparaison : comme en JS, mais en pire. TOUJOURS === et !==.
var_dump(0 == 'a');           // PHP 8 : false (avant PHP 8 : TRUE -- un CVE en puissance)
var_dump('12abc' == 12);      // false en PHP 8 (true en PHP 7)
var_dump(100 === '100');      // false
var_dump(NAN === NAN);        // false
var_dump(0.1 + 0.2 === 0.3);  // false ; compare avec abs(a-b) < PHP_FLOAT_EPSILON
\end{codephp}

\subsection{Cha\^ines : tout ce qu'on utilise vraiment}

\begin{codephp}
<?php
$s = "  Le Chat et le CHIEN  ";

strlen($s);                       // en OCTETS ("\'e" = 2 octets en UTF-8) -> ne l'utilise pas pour l'affichage
mb_strlen($s, 'UTF-8');           // en caracteres reels ; mb_* = l'extension indispensable
mb_strtolower($s);  mb_convert_case($s, MB_CASE_TITLE, 'UTF-8');
str_contains($s, 'chat');          // 8.0 : remplace strpos !== false
str_starts_with($s, '  Le');  str_ends_with($s, ' ');
trim($s);  ltrim($s);  rtrim($s, " \t\n\r\0\x0B");
str_replace(['le', 'et'], ['LES', 'ET'], strtolower($s));
str_repeat('-', 10);  substr($s, 2, 5);  substr_count($s, 'e');
strpos(strtolower($s), 'chat');     // position ou false
explode(',', 'a,b,c');  implode(', ', ['a', 'b']);
sprintf('%s : %05d | %-10s | %+.2f', 'facture', 42, 'total', -12.5);
number_format(1234.5678, 2, ',', ' ');        // "1 234,57" (fr)
nl2br($texteUtilisateur);                       // sauts de ligne -> <br>
strip_tags($html);  htmlspecialchars($html, ENT_QUOTES | ENT_SUBSTITUTE, 'UTF-8');   // echapper pour le HTML
wordwrap($s, 80, "\n", true);  str_pad('7', 3, '0', STR_PAD_LEFT);
str_word_count($s);  ucfirst($s);  ucfirst(mb_strtolower($s));

// Interpolation et formatage multi-lignes
$prenom = 'Nour';
echo "Bonjour $prenom, tu as " . strlen($prenom) . " lettres.\n";
echo "Detail : {$article['titre']} | {$user->nom}\n";
$sql = <<<SQL
    SELECT id, titre
    FROM articles
    WHERE publie = 1
    ORDER BY id DESC
    SQL;                     // 7.3+ : le cloturant peut etre indentes
$md = <<<'TXT'               // nowdoc = comme une chaine simple, aucune interpolation
Texte litteral : $variable, {$choses}, \n
TXT;
\end{codephp}

\begin{piege}
\textbf{Encodage.} PHP manipule des octets, pas des caract\`eres. Sans \verb|mbstring| et sans
UTF-8 coh\'erent, tu auras des \og ~\~A~ \fg{} partout. La checklist~: fichiers en UTF-8 sans BOM~;
base en \verb|utf8mb4| avec \verb|utf8mb4\_unicode\_ci|~; en-t\^ete \verb|Content-Type: text/html; charset=utf-8|~;
et les fonctions \verb|mb\_*| pour tout ce qui touche au texte humain.
\end{piege}

\section{Tableaux~: le vrai type universel de PHP}

Un tableau PHP est une \textbf{table ordonn\'ee}~: il accepte des cl\'es (entiers ou cha\^ines) et
maintient l'ordre d'insertion. C'est \`a la fois un tableau, une liste, un dictionnaire et une file.

\begin{codephp}
<?php
$article = [
    'id' => 7,
    'titre' => 'Apprendre PHP',
    'tags' => ['php', 'web'],
    'auteur' => ['nom' => 'Nour', 'email' => 'nour@example.com'],
    0 => 'premiere cle numerique -- attention, "0" devient 0',
];

// Acces et ecriture
$article['titre'];
$article['auteur']['nom'];
$article['vue'] ??= 0;                    // 7.4+ : initialise si absent/nul
$article['titre'] .= ' (v2)';
unset($article['tags']);

// Creation et fusion
$nouveau = array_fill_keys(['a', 'b'], 0);
$combine = array_combine(['id', 'nom'], [1, 'Nour']);
$fusion = $article + ['id' => 999];                 // les cles existantes gagnent
$fusion2 = array_merge($a, $b);                    // reindexe les cles numeriques  <-- piege
$deep = array_replace_recursive($config, $surcharges);
$tabs = [...$premiers, ...$derniers];              // spread : seulement listes (cles numeriques seq.)

// Cles et valeurs
array_keys($article);  array_values($article);
count($article);  in_array('php', $article['tags'] ?? [], true);   // le 3e parametre = comparaison stricte
array_key_exists('id', $article);     // != isset() : lui voit les cles a valeur null

// Transformer (les trois qui remplacent 90 % des boucles)
$mots = explode(' ', strtolower(trim($s)));
$longueurs = array_map(mb_strlen(...), $mots);
$filtr\'es = array_filter($mots, static fn (string $m): bool => mb_strlen($m) > 3);
$reindexe = array_column($articles, null, 'id');           // clef = id -> recherche O(1)
$titre = array_map(static fn (array $a): string => $a['titre'], $articles);
$total = array_reduce($lignes, static fn (float $c, array $l): float => $c + $l['prix'] * $l['qte'], 0.0);

// Ordre
sort($liste);                        // par valeur, reindexe -- MUTANT
asort($tableauAssoc);                // par valeur, garde les cles
ksort($parClef);
usort($articles, static fn (array $a, array $b): int => $b['date'] <=> $a['date']);
// spaceship <=> renvoie -1/0/1 : la seule bonne facon de comparator en PHP
array_multisort(array_column($articles, 'note'), SORT_DESC, $articles);   // multi-criteres

// Decoupage, piles, files
array_slice($liste, 2, 5, preserve: true);
array_chunk($articles, 20);
array_splice($liste, 1, 2, $remplacement);
array_push / array_pop / array_shift / array_unshift ;
$array = [1, 2, 3]; $x = array_pop($array);      // 3
end($array) / reset($array) / next($array) / current($array) / key($array);
list($a, $b) = [1, 2];   [$x, $y] = [1, 2];    ['nom' => $n, 'id' => $i] = $article;   // destructuration
\end{codephp}

\section{Conditions, boucles, et le flux de contr\^ole}

\begin{codephp}
<?php
$note = 14;
if ($note >= 16) {
    echo 'Tres bien';
} elseif ($note >= 10) {
    echo 'Admis';
} else {
    echo 'Releve';
}

// alternative syntax : tres lisible dans les vues (HTML mele au PHP)
foreach ($articles as $i => $article): ?>
    <article class="<?= $article['publie'] ? 'publie' : 'brouillon' ?>">
        <h2><?= htmlspecialchars($article['titre']) ?></h2>
        <p><?= e($article['extrait'] ?? '') /* fonction utilitaire a definir */ ?></p>
    </article>
<?php endforeach;

switch ($etat) {                    // switch traditionnel : comparaison loose (==) -- evite-le
    case 'ok': break;
    default: throw new InvalidArgumentException("etat inconnu : {$etat}");
}
$resultat = match ($note) {        // 8.0 : comparaison stricte, expression, pas de fallthrough
    20 => 'parfait',
    16, 17, 18, 19 => 'tres bien',
    10 ... 15 => 'bien',            // plage
    default => 'a revoir',
};
$classe = match (true) {
    $note >= 16 => 'excellent',
    $note >= 10 => 'moyen',
    default => 'faible',
};

for ($i = 0; $i < 10; $i++) {}
while ($ligne = fgets($h)) {}        // idiomatique : affectation + test
do { $n++; } while ($n < 10);
foreach ($tableau as $clef => $valeur) {}
foreach ($generateur() as $ligne) {}    // memoire constante sur un gros fichier
for (;;) { break; }                     // boucle infinie propre
\end{codephp}

\section{Fonctions, fermetures, appelables}

\begin{codephp}
<?php
declare(strict_types=1);

// Fonction nommee, typ\'ee, avec valeurs par defaut
function envoyerEmail(
    string $destinataire,
    string $sujet,
    string $corps = '',
    array $piecesJointes = [],
    ?DateTimeImmutable $programmePour = null,
): bool {
    if (!filter_var($destinataire, FILTER_VALIDATE_EMAIL)) {
        throw new InvalidArgumentException("email invalide : {$destinataire}");
    }
    return true;    // ici, l'appel a la fonction mail() du serveur SMTP reel
}

// Nomme (8.0) : l'API devient lisible, l'ordre des parametres ne compte plus
envoyerEmail(destinataire: 'nour@example.com', piecesJointes: ['facture.pdf'], sujet: 'Facture');

// Passage par valeur (defaut) vs par reference (a eviter, sauf cas precis)
function ajout(int &$compteur): void { $compteur++; }
function tableauEnReferenc(array &$t): void { $t['vu'] = true; }

// Fermetures : "use" importe le contexte (par copie ; par reference avec "&")
$tva = 0.2;
$ht = [10, 20, 30];
$ttc = array_map(static fn (float $x): float => $x * (1 + $tva), $ht);
$compteur = function () use (&$n): int { return ++$n; };     // le "&" est essentiel ici
$appliquer = fn (callable $f, $v) => $f($v);                  // fleche : capture auto, une expression

// Generateur : produire une sequence sans la stocker
function lignes(string $chemin): Generator
{
    $h = fopen($chemin, 'rb');
    if ($h === false) { throw new RuntimeException("lecture impossible : {$chemin}"); }
    try {
        while (($l = fgets($h)) !== false) {
            yield rtrim($l, "\r\n");
        }
    } finally {
        fclose($h);
    }
}
foreach (lignes('acces.log') as $n => $ligne) { /* memoire = 1 ligne */ }

// Callables : le PHP n'aime pas les fonctions de haut niveau silencieuses
$tri = usort($articles, 'comparerParDate');            // nom de fonction
$tri = usort($articles, [$objet, 'methode']);
$tri = usort($articles, [MaClasse::class, 'statique']);
$tri = usort($articles, $objet->methode(...));          // premier-class callable (8.1) : le plus propre
$tri = usort($articles, fn (array $a, array $b) => $b['id'] <=> $a['id']);

// Functions utiles a connaitre : call_user_func est depasse par $cb(...) ; array_* le reste.
\end{codephp}

\subsection{Les 15 fonctions globales qu'on cherche toujours}

\begin{tabularx}{\textwidth}{@{}l X@{}}
  \toprule
  \textbf{Fonction} \& \textbf{Usage}\\
  \midrule
  \verb|count|, \verb|iterator\_to\_array|, \verb|is\_array| \& dimensions et v\'erifications.\\
  \verb|array\_map|, \verb|array\_filter|, \verb|array\_values| \& transformations.\\
  \verb|str\_contains|, \verb|str\_replace|, \verb|sprintf|, \verb|trim| \& cha\^ines.\\
  \verb|htmlspecialchars| (\verb|ENT\_QUOTES|) \& sortir du texte en HTML.\\
  \verb|json\_encode| / \verb|json\_decode|, avec \verb|JSON\_THROW\_ON\_ERROR| \& API et configs.\\
  \verb|filter\_var| / \verb|filter\_input| \& validation d'email, d'URL, d'entier.\\
  \verb|password\_hash| / \verb|password\_verify| \& mots de passe.\\
  \verb|random\_int| / \verb|random\_bytes| \& al\'eatoire \emph{cryptographique} (jamais \verb|rand| pour un jeton).\\
  \verb|file\_get\_contents| / \verb|file\_put\_contents| / \verb|is\_readable| \& fichiers simples.\\
  \verb|date|, \verb|DateTimeImmutable| \& dates~; pr\'ef\'erence absolue \`a la classe.\\
  \verb|var\_dump|, \verb|print\_r|, \verb|error\_log| \& debug (et \verb|dump()| de Symfony).\\
  \verb|compact|, \verb|extract| \& variables vers une vue (le second est \`a fuir).\\
  \verb|enum|? \verb|enum\_exists|, \verb|class\_exists|, \verb|function\_exists| \& introspection.\\
  \verb|match|, \verb|throw expr| (8.0) \& flux de contr\^ole.\\
\end{tabularx}

\section{TP --- un mini-routage et un rapport CSV}

Deux fichiers, sans base de donn\'ees, sans framework. L'id\'ee~: voir que PHP n'a besoin de rien
de plus pour faire un vrai site.

\begin{codephp}
<?php
// public/index.php --- le "front controller" : TOUT passe par la.
declare(strict_types=1);
error_reporting(E_ALL);
ini_set('display_errors', '1');          // en DEV seulement ; en prod : 0 + logs

session_start();
$uri = parse_url($_SERVER['REQUEST_URI'], PHP_URL_PATH) ?: '/';
$methode = $_SERVER['REQUEST_METHOD'];

// Une petite fonction d'echappement, reutilisee dans les vues
function e(?string $v): string { return htmlspecialchars((string) $v, ENT_QUOTES | ENT_SUBSTITUTE, 'UTF-8'); }

$routes = [
    'GET /' => 'accueil',
    'GET /articles' => 'articles',
    'POST /articles' => 'creerArticle',
    'GET /rapport.csv' => 'rapport',
];
$action = $routes["$methode $uri"] ?? null;
if ($action === null) {
    http_response_code(404);
    require __DIR__ . '/../src/vues/404.php';
    exit;
}
require __DIR__ . "/../src/actions/{$action}.php";
\end{codephp}

\begin{codephp}
<?php
// src/actions/articles.php
declare(strict_types=1);
$cheminDonnees = __DIR__ . '/../../data/articles.json';
$articles = json_decode((string) file_get_contents($cheminDonnees), true, 512, JSON_THROW_ON_ERROR);
$filtre = trim((string) ($_GET['q'] ?? ''));
$parPage = 10;
$page = max(1, (int) ($_GET['page'] ?? 1));

$q = array_values(array_filter(
    $articles,
    static fn (array $a): bool => $filtre === ''
        || mb_stripos($a['titre'], $filtre) !== false
        || mb_stripos($a['contenu'] ?? '', $filtre) !== false
));
$total = count($q);
$nbPages = max(1, (int) ceil($total / $parPage));
$page = min($page, $nbPages);
$articles = array_slice($q, ($page - 1) * $parPage, $parPage);

require __DIR__ . '/../vues/liste.php';
\end{codephp}

\begin{codephp}
<?php
// src/actions/rapport.php --- exporter un CSV correctement (le BOM UTF-8 est un detail reel)
declare(strict_types=1);
$articles = json_decode((string) file_get_contents(__DIR__ . '/../../data/articles.json'), true, 512, JSON_THROW_ON_ERROR);

$entetes = ['id', 'titre', 'auteur', 'date', 'vues'];
$corps = fopen('php://output', 'wb');
fwrite($corps, "\xEF\xBB\xBF");                     // BOM : Excel ouvre alors en UTF-8
fputcsv($corps, $entetes, ',', '"', '\\');           // 1ere colonne forcee en chaine, 4e argument = 8.4
foreach ($articles as $a) {
    fputcsv($corps, [
        (string) $a['id'],
        $a['titre'],
        $a['auteur']['nom'] ?? '',
        date('Y-m-d', (int) ($a['cree'] ?? time())),
        (int) ($a['vues'] ?? 0),
    ]);
}
fclose($corps);
exit;
\end{codephp}

\begin{note}
Les \verb|require|/ \verb|include| ont deux variantes. \verb|require| \'echoue fatalement,
\verb|include| \'emet un avertissement~; \verb|*\_once| \'evite le double-chargement. \verb|\_\_DIR\_\_| rend
le chemin insensible au r\'epertoire courant~: utilise-le syst\'ematiquement, sinon ton application
casse d\`es que le serveur est lanc\'e depuis un autre dossier.
\end{note}

\subsection{Les erreurs de PHP que les d\'ebutants rencontrent}

\begin{tabz}{lX}
  \headrow \thd{Message} \& \thd{Cause r\'eelle} \\
  \midrule
  \verb|Cannot modify header information - headers already sent| \& Un espace, un retour \`a la ligne ou un \verb|echo| avant \verb|header()|~: v\'erifie l'absence de \verb|?>| final et d'octets parasites avant \verb|<?php|. \\
  \verb|undefined array key "x"| \& la cl\'e n'existe pas~; utilise \verb|??| ou \verb|array\_key\_exists|. \\
  \verb|TypeError: Return value must be of type int, string returned| \& \verb|strict\_types=1| fait son travail~: convertis explicitement avec \verb|(int)| ou \verb|intcap|. \\
  \verb|Allowed memory size exhausted| \& boucle infinie, tableau gigantesque, ou r\'ecursion sans fin. \\
  Page blanche \& une erreur fatale avec \verb|display\_errors=Off|. Regarde le journal~: \verb|tail -f /var/log/php-fpm/error.log|, ou active \verb|error\_reporting(E\_ALL)| + \verb|display\_errors=1|. \\
  \verb|Class "X" not found| \& autoload non configur\'e (Composer n'est pas install\'e/charg\'e), ou nom d'espace de lignes incorrect. \\
\end{tabz}

\begin{exercice}
\begin{enumerate}
  \item \'Ecris \verb|slugify(string \$t): string|~: minuscules, accents retir\'es
        (\verb|iconv('UTF-8', 'ASCII//TRANSLIT//IGNORE', \$t)| ou \verb|mb\_convert\_encoding|),
        non-alphanum\'eriques en tirets, tirets dupliqu\'es et marges nettoy\'es~; test\'e sur
        \og ~L'\'Ev\'enement de l'ann\'ee, \`a Thi\`es !~ \fg{}.
  \item \'Ecris \verb|formaterDuree(int \$secondes): string| renvoyant \og ~2 h 05 min~ \fg{}, et
        \verb|nombreEnLettres(int \$n): string| pour 0--999 (r\'ecursif).
  \item Construis \verb|tableauVersCsv(array \$lignes): string| et \verb|csvVersTableau(string): array|
        qui se comportent correctement avec guillemets, virgules et retours \`a la ligne dans les champs.
  \item Analyse un fichier \verb|access.log|~: renvoie les 10 IP les plus fr\'equentes, le total de
        requ\^etes par code HTTP, et la liste des 404, avec un g\'en\'erateur pour tenir sur un fichier
        de 1~Go.
  \item \'Ecris une fonction \verb|validerFormulaire(array \$donnees, array \$regles): array|~: r\`egles
        \verb|required|, \verb|min:5|, \verb|email|, \verb|numeric|, \verb|in:a,b,c|,
        \verb|confirmed|~; retour \verb|[champ => [erreurs]]| vide si tout passe. C'est la base
        de ton futur micro-framework.
  \item \'Ecris le \verb|creerArticle.php| de l'exemple~: v\'erifie la m\'ethode, valide, ajoute au JSON
        avec un verrou (\verb|flock|) et une \'ecriture atomique (\verb|tempnam| + \verb|rename|),
        renvoie une redirection (303) vers la page de l'article.
\end{enumerate}
\end{exercice}

\begin{ressources}
\begin{itemize}
  \item Le manuel officiel (avec les notes d'utilisateurs~; la 2\textsuperscript{e} meilleure doc du web) --- \url{https://www.php.net/manual/fr/}
  \item PHP The Right Way (traduit, \`a lire en entier~: les bonnes pratiques modernes) --- \url{https://phptherightway.com/}
  \item PHPFIG standards --- \url{https://www.php-fig.org/psr/}
  \item php-pdo / PDO---\url{https://www.php.net/manual/en/book.pdo.php}
  \item \link{PHP~Watched}{https://php.watch/} (nouveautes par version) et \url{https://stitcher.io/blog/php-84-in-a-hurry}
\end{itemize}
\end{ressources}

\index{PHP!bases}\index{array}\index{strict\_types}
